% MAIN SNAPSHOT — LSPV/RMV::VALSUBREVIEW1::2026-10-14T00:48:11+09:00
% !TEX program = xelatex
\documentclass[11pt]{article}
\usepackage[T1]{fontenc}
\usepackage{lmodern,amsmath,amssymb,amsthm,booktabs,tabularx,xurl,hyperref}
\usepackage[margin=0.86in]{geometry}
\newtheorem{theorem}{Theorem}
\newtheorem{corollary}[theorem]{Corollary}
\newcommand{\F}{\mathbb F_2}
\newcommand{\EXACT}{\textbf{PROVED/EXACT}}
\newcommand{\KNOWN}{\textbf{KNOWN THEOREM}}
\newcommand{\COMP}{\textbf{COMPUTATIONAL}}
\newcommand{\COND}{\textbf{CONDITIONAL}}
\newcommand{\CAND}{\textbf{CANDIDATE}}
\newcommand{\OPEN}{\textbf{OPEN}}
\newcommand{\RET}{\textbf{FALSIFIED/RETIRED}}
\hypersetup{colorlinks=true,urlcolor=blue,linkcolor=blue}
\setlength{\emergencystretch}{3em}
\title{VALSUBREVIEW1: Fixed-Valuation Cube Selector Review\\
and a Logarithmic-Subfield Extension}
\author{UNIVERSAL\_CORE\_SYNTHESIS / LSPV-5}
\date{14 October 2026, 00:48 KST}

\begin{document}
\maketitle

\begin{center}
\fbox{\parbox{0.95\linewidth}{
\textbf{Executive status: P versus NP remains OPEN.}
Independent synthesis review validates VALSUBCUBEFW1's fixed-$\nu_3(d)$
selector, promoted here to \EXACT{}.  The argument also yields a stronger
uniform-window theorem: when
$3^{\nu_3(d)}=O(\log d)$, typed promised-image cube roots in
$\mathbb F_{2^d}$ are still deterministically selectable in polynomial time.
That extension has a complete proof here but is retained as \CAND{} pending
independent specialist review.  Thus neither $6\mid d$ nor slowly unbounded
$3$-adic valuation repairs the residual lower-bound route.  No unrestricted
all-program lower bound follows.}}
\end{center}

\section{Formal claim card}
\begin{tabularx}{\linewidth}{p{0.27\linewidth}X}
\toprule
Claim ID & VALSUBREVIEW1 \\
Status before run & Fixed-$\nu_3(d)$ selector \CAND; unrestricted cube-root and machine lower bounds \OPEN \\
Desired conclusion & Review the fixed-valuation proof and identify the largest immediate uniform polynomial escape \\
Exact quantifiers & (i) For every fixed $a\ge1$, one uniform deterministic program works for every even $d$ with $\nu_3(d)=a$ and every typed degree-$d$ field; (ii) for every fixed $C>0$, one uniform program works for all even $d$ satisfying $3^{\nu_3(d)}\le C\log_2(d+2)$ \\
Input family & Explicit irreducible polynomial-basis encodings of $\mathbb F_{2^d}$, target $b$, and the already-declared typed Tseitin/cotree interface \\
Machine/model & Uniform deterministic random-access bit computation with polynomial workspace and full input access \\
Resource & Parsing, valuation, Frobenius matrix, Gaussian elimination, subfield enumeration, exponentiation, correction, image verification, forest reconstruction, and output \\
SAME\_REALIZATION & The same typed field, target, selected root, charge vector, cotree assignment, and reconstructed edge witness \\
Falsifier & A legal degree in the stated window for which the constructed subfield lacks the full $3$-primary order, or a promised cube for which all correction candidates fail \\
\bottomrule
\end{tabularx}

\section{Primary literature and transfer audit}
Adleman--Manders--Miller introduced finite-field root extraction with a
supplied nonresidue \cite{AMM}; Cao--Sha--Fan spell out the general
$r$-th-root correction and its discrete-log component \cite{Cao}.  Sze gives
deterministic root extraction without a supplied nonresidue under explicit
factor and root-of-unity hypotheses, while warning that efficiency depends on
that arithmetic regime \cite{Sze}.  Bhargava--Ivanyos--Mittal--Saxena relate
nonresidue construction, extension construction, and deterministic root
finding in other parameter regimes \cite{BIMS}.

The first two mechanisms are \textbf{directly transferable}: the residual is
literally a power equation in the represented field.  Sze is a
\textbf{useful structural analogue requiring parameter verification}; this
report does not import a broader theorem from it.  The specific BIMS
$r\mid\gcd(m,p-1)$ root-finding corollary is a \textbf{failed direct
transfer} to $p=2,r=3$, although its nonresidue viewpoint is relevant.
Cryptographic cube-root hardness is a \textbf{misleading analogy}: the
explicit subfield is exploitable structure, not an admissible assumption.

\section{Independent review of the fixed-valuation theorem}
Let $a\ge1$ be fixed, $d$ be even with $\nu_3(d)=a$, and
$K=\mathbb F_{2^d}$ in the supplied polynomial basis.  Put
\[
 r=2\cdot3^a,\qquad N=2^d-1=3^{a+1}m,\qquad 3\nmid m.
\]
Because $d/3^a$ is even, $r\mid d$.  Lifting the exponent gives
\[
 \nu_3(2^d-1)=1+\nu_3(d)=a+1=\nu_3(2^r-1).             \tag{1}
\]
The kernel of the $\F$-linear map $z\mapsto z^{2^r}+z$ is therefore the
unique subfield $K_0\cong\mathbb F_{2^r}$.  Its multiplicative cofactor
\[
 S=\frac{2^d-1}{2^r-1}
\]
is prime to $3$ by (1).  Hence a generator of $K_0^*$ is not a cube in $K$.
Enumerating $K_0$ and testing $c^{N/3}\ne1$ deterministically finds such a
$c$.  For fixed $a$, this enumeration has family-dependent constant size.

Set $\eta=c^m$.  It has order exactly $3^{a+1}$.  For a nonzero target $b$,
let $e=3^{-1}\bmod m$ and $y=b^e$.  If $b$ is a cube, then
$y^3/b$ belongs to the unique subgroup $\langle\eta^3\rangle$ of order
$3^a$.  Consequently exactly one member of the transversal
\[
 x_j=y\eta^j,\qquad 0\le j<3^a,                         \tag{2}
\]
satisfies $x_j^3=b$.  Direct recomputation accepts that member and rejects an
off-image target; zero maps to zero.  Constructing the Frobenius matrix,
finding its kernel, enumerating the fixed subfield, exponentiating, testing
(2), and writing the result are polynomial in $d$ for fixed $a$.

Every algebraic and complexity step in VALSUBCUBEFW1 is therefore correct.
Its fixed-$\nu_3(d)$ selector is promoted to \EXACT{} after this independent
synthesis review.  This promotion concerns a root-selection upper bound, not
SAT source search and not a separation.

\section{A stronger uniform logarithmic-subfield window}
\begin{theorem}[logarithmic-subfield selector]
Fix $C>0$.  There is a uniform deterministic algorithm which, for every even
$d\ge4$ satisfying
\[
 3^{\nu_3(d)}\le C\log_2(d+2),                          \tag{3}
\]
every explicit irreducible degree-$d$ polynomial, and every
$b\in\mathbb F_{2^d}$, returns an $x$ with $x^3=b$ iff $b$ is a cube and
otherwise rejects.  Its bit runtime and workspace are polynomial in the
typed input length, with the polynomial exponent allowed to depend on $C$.
\end{theorem}
\begin{proof}
Compute $a=\nu_3(d)$.  If $a=0$, use the verified image-subgroup exponent
selector of CUBEROOTSEL1.  If $a\ge1$, run the construction above with
$r=2\cdot3^a$.  This is now one input-driven algorithm rather than a separate
program for each $a$.  The only formerly nonuniform-looking cost is
enumeration of $K_0$, but (3) gives
\[
 2^r=2^{2\cdot3^a}\le (d+2)^{2C}.                      \tag{4}
\]
The correction search has $3^a\le C\log_2(d+2)$ candidates.  Frobenius
evaluation, $d\times d$ linear algebra, polynomial-basis field operations,
image verification, and output are polynomial in $d$.  Multiplying those
costs by (4) remains $d^{O_C(1)}$.  Equations (1)--(2) prove correctness, and
the final cube check makes off-image behavior total.
\end{proof}

The proof is complete, but because this extension is new in this synthesis
cycle it remains \CAND{} pending an independent specialist review.  It
strictly enlarges the firewall: $\nu_3(d)$ may tend to infinity as slowly as
$O(\log\log d)$ subject to (3).  Failure of (3) is only failure of this
enumeration bound; it is not a lower bound for the complementary degrees.

\section{Bounded reproducible audit}
The independent script
\path{experiments/review/fixed_valuation_selector_synthesis_audit.py}
ran once in 0.022 seconds with seed 14102026.  It checks the cyclic-exponent
identities independently of the predecessor's polynomial-basis and Tseitin
implementation.  It passed 16,258 assertions with zero failures:
\begin{itemize}
\item 50 even degrees through 300 with $a=1,2,3,4$;
\item 7,230 promised cubes, 1,200 off-image exponents, and 4,878 selected
roots different from the sampled source root;
\item exact subfield-order, cofactor, noncube, $\eta$-order, and unique
correction-transversal checks, with up to 81 candidates;
\item 66 genuine square-cotree valuation checks; and
\item eight symbolic-size controls with unbounded valuations $a=1,\ldots,8$
and $r=O(\log d)$.
\end{itemize}
The audit is \COMP{}.  It does not prove either asymptotic theorem; the
displayed proofs do.  Raw output is
\path{results/raw/review/2026-10-14_0048_fixed_valuation_selector_synthesis.json};
the code SHA-256 is
\texttt{f1fc5430fe3a6a9ffdab2b7fb4391561ec03307441ddf422585426fe98aa93c2}.
The predecessor's separate exact-field audit passed 7,659 assertions,
including complete $K_{7,7}$ witness lifts.

\section{Barrier, wrappers, and post-quantum ledger}
The selector allows full source and target access, arbitrary polynomial
workspace, compact field arithmetic, and charged discovery, reconstruction,
verification, and output.  DEAD/COPY/DELAY wrappers do not remove it;
INTERP/RECOMP can execute it.  A semantic or algebraic oracle can only make
the upper bound easier, so this firewall relativizes and algebrizes.  Natural
Proofs is not invoked.

For the post-quantum ledger, finite-field multiplication, Frobenius-fixed
subfields, and their $3$-primary relation are an \textbf{EXACT
OBJECT/THEOREM TRANSFER}.  Schur-square, quadratic-relation, Hilbert/syzygy,
collision-entropy, and LWE-style extraction language is a \textbf{FAILED
TRANSFER} to unrestricted runtime here: the same hidden algebraic structure
is algorithmically exploitable.  Generic binding is only an
\textbf{ANALOGY/CANDIDATE} without a construction-specific forbidden witness.

\section{CORPUS REVIEWED}
This run surveyed the current GitHub tree; root \path{main.tex};
\path{AGENTS.md}, \path{CODEX.md}, the complete PNP research skill and its
mode/protocol references; all seven task specifications; newest MAIN
PROMISEGATE1 and historical CUBEPERMFW1; AFFINE-RESIDUAL-REDIRECT1,
EARLY22-INTEGRATION1, SINGLE22-SEPARATION1, and AFFINE-RESIDUAL-RANK1; and the
HEC/VSC-HEC/G-PPC-HEC, CRS-PPC, RCCN, CCSPC, and Search-G transfer guards.
The actual current specialist evidence read was CHECKABLEPERM1,
PLANTEDCUBEFW1, CUBEROOTSEL1, and VALSUBCUBEFW1 with cited code/raw output.
The exact-$(2,2)$ compilers and proof projections remain preserved only in
their declared proof systems.

Predecessor recovery was not required: CHECKABLEPERM1 exists in LSPV-4's
canonical directory, is substantive, and contains the required header and
final sentinel.  Its strict recognizer repair is retained; this cycle does
not independently promote its review-pending semantic claims.

\section{DUPLICATION / CONTRADICTION AUDIT}
Repository searches for fixed valuation, subfield cubic nonresidues, AMM
correction, and unbounded valuation found no earlier project theorem beyond
VALSUBCUBEFW1.  CUBEROOTSEL1 treats $a=0$; PLANTEDCUBEFW1 uses a source
section; CUBEPERMFW1 treats odd $d$.  The uniform window (3) is new to the
corpus.  It contradicts no proof-system lower bound because it is a witness
selector for a special typed residual, not an arbitrary proof simulation.

\section{WHAT CHANGED RELATIVE TO ROOT MAIN}
\begin{itemize}
\item The fixed-$\nu_3(d)$ subfield selector is promoted from \CAND{} to
\EXACT{} after a second proof review.
\item The residual firewall now includes slowly unbounded valuations obeying
(3); that extension is \CAND{} pending independent review.
\item ``Unbounded $\nu_3(d)$'' is too coarse as the next open regime.  Only
the super-logarithmic-subfield range $3^{\nu_3(d)}/\log d\to\infty$ escapes
this particular enumeration argument.
\item The unrestricted lower-bound arrow and P-versus-NP status are unchanged.
\end{itemize}

\section{ADOPTED / REJECTED / RETIRED PRIOR COMPONENTS}
\textbf{Adopted:} strict typed field/graph interfaces, the exact forest lift,
CUBEROOTSEL1, the fixed-subfield derivation, and full cost accounting.
\textbf{Preserved:} EARLY22/SINGLE22 compiler and PC/PCR/Resolution results
at their restricted scope, plus the conditional growing-factor lemma.
\textbf{Rejected:} uncharged nonresidue advice, source-preimage advice,
cryptographic hardness, and failure of one monomial inverse as evidence of
hardness.  \textbf{Retired:} $6\mid d$, fixed valuation, and slowly growing
valuation satisfying (3) as sufficient cube-residual repairs.

\section{STALE OR OVERSTATED CLAIMS FOUND}
Root and CUBEROOTSEL1 describe the ``remaining $6\mid d$ case'' too broadly.
VALSUBCUBEFW1 already removes every fixed valuation, and theorem (3) removes
part of the unbounded-valuation range.  The honest remaining algorithmic
question is not all $6\mid d$ but the complementary degrees where
$3^{\nu_3(d)}$ grows faster than every allowed logarithmic bound, together
with alternative nonresidue-free algorithms.  Even there, lack of a known
selector would not prove hardness.

\section{Current short proof chain and status}
\begin{center}
\begin{tabularx}{0.98\linewidth}{p{0.40\linewidth}p{0.19\linewidth}X}
\toprule
Claim & Status & Scope \\
\midrule
Affine elimination and forest witness & \EXACT & Every consistent explicit Tseitin charge \\
Projected rank and $2^{\rho+h}$ escape & \EXACT & Declared support and external variables \\
CUBEROOTSEL1, $\nu_3(d)=0$ & \EXACT & Independent LSPV-3 review complete \\
Fixed-$\nu_3(d)$ subfield selector & \EXACT & Independent synthesis review complete \\
Logarithmic-subfield selector (3) & \CAND & Complete proof; independent review pending \\
Finite synthesis audit & \COMP & 16,258 assertions, zero failures \\
Growing-factor implication & \COND & Sufficient premise remains uninstantiated \\
Unrestricted all-program lower bound & \OPEN & No model-matched theorem \\
P versus NP separation & \OPEN & Load-bearing implication remains open \\
\bottomrule
\end{tabularx}
\end{center}

The shortest honest chain is
\[
 B_Xx=c\quad\Longrightarrow\quad x=a+Lz
 \quad\Longrightarrow\quad Q(u,v,a+Lz),
\]
followed, on the cube subfamilies above, by a charged deterministic selector
and forest reconstruction.  This is an upper-bound firewall.  It supplies no
net multiplier for $H$ and no contradiction to $P=NP$.

\textbf{Single highest-value next assignment.}
Independently review theorem (3), then attack the complementary
super-logarithmic-subfield degrees with Sze-style nonresidue-free extraction
and explicit $3$-power subfield towers.  Either construct another uniform
selector or isolate the precise deterministic algorithmic obstruction.  Do
not promote failure of known root algorithms to an unrestricted lower bound.

\begin{thebibliography}{9}
\bibitem{AMM} L. M. Adleman, K. L. Manders, and G. L. Miller,
``On Taking Roots in Finite Fields,'' \emph{18th FOCS}, 1977, 175--178.
\bibitem{Cao} Z. Cao, Q. Sha, and X. Fan,
``Adleman--Manders--Miller Root Extraction Method Revisited,''
arXiv:1111.4877, 2011.
\bibitem{Sze} T.-W. Sze,
``On Taking $r$-th Roots without $r$-th Nonresidues over Finite Fields and
Its Applications,'' arXiv:1105.5852, 2011.
\bibitem{BIMS} V. Bhargava, G. Ivanyos, R. Mittal, and N. Saxena,
``Irreducibility and $r$-th Root Finding over Finite Fields,''
arXiv:1702.00558, 2017.
\end{thebibliography}

\bigskip
\noindent P versus NP remains open.

\end{document}
% END_OF_VERIFIED_MAIN::VALSUBREVIEW1
